\chapter{概率图模型概述}
\label{chap:pgm-overview}

一台设备发出报警，是否应当立即停机？报警可能来自设备故障造成的过热，也可能来自传感器失灵。只输出“有故障”或“无故障”，无法说明另一个结果还有多大可能，也无法回答新增证据会怎样改变判断。即使分类结果相同，故障概率为$0.02$和$0.45$时，维修人员面对的风险也很不同。

要回答这类问题，需要描述多个不确定因素怎样共同变化。\term{概率图模型}（probabilistic graphical model, PGM）用图组织随机变量之间的局部依赖，再用局部概率函数组成联合分布。图使模型不必为所有变量组合分别指定概率，也让后续计算能够利用局部结构。本章只建立后续章节共同需要的概念：为什么要建模分布，图模型怎样表示联合分布，以及推断和学习分别解决什么问题。

\section{从单值预测到结构化分布}
\label{sec:pgm-overview-distribution-prediction}

熟悉的预测形式是$\hat y=f(x)$：给定输入$x$，模型返回一个数值或类别。例如，以传感器读数为输入，输出设备下一时刻的温度。这种\term{点预测}（point prediction）把结果压缩成一个值，适合直接给出答案，却不保留其他结果的可能性。

\term{预测分布}（predictive distribution）$p(y\mid x)$描述输入固定后，不同输出有多大可能。对于离散输出，它为每种结果分配概率；对于连续输出，它给出概率密度，区间概率由积分得到。分布不仅告诉我们结果集中在哪里，还能表示多个彼此分离的高概率区域。

图~\ref{fig:pgm-overview-predictive-density}给出一个人工构造的温度预测。冷却系统可能处于两种工况，预测密度分别集中在$40\,^\circ\mathrm{C}$和$80\,^\circ\mathrm{C}$附近，权重各为一半。均值$60\,^\circ\mathrm{C}$位于两个峰之间，并不是高概率区域。若只保留均值，就看不见高温工况；若只保留一个峰，又会丢掉另一种工况。这里的分布为教学设定，不是实际设备的测量结果。

\begin{figure}[htbp]
  \centering
  \begin{tikzpicture}
  \begin{axis}[
    width=105mm,
    height=58mm,
    xmin=25,xmax=95,
    ymin=0,ymax=0.08,
    axis lines=left,
    axis line style={BookMuted},
    tick style={BookMuted},
    tick label style={font=\small,text=BookMuted},
    label style={font=\small,text=BookInk},
    xlabel={温度$y$（$^\circ\mathrm{C}$）},
    ylabel={概率密度（$^\circ\mathrm{C}^{-1}$）},
    xtick={40,60,80},
    ytick={0,0.04,0.08},
    scaled ticks=false,
    clip=false
  ]
    \addplot[
      draw=BookTeal,
      line width=1pt,
      domain=25:95,
      samples=181
    ] {0.5/(3*sqrt(2*pi))*(exp(-(x-40)^2/18)+exp(-(x-80)^2/18))};
    \draw[BookAmber,dashed,line width=0.8pt]
      (axis cs:60,0) -- (axis cs:60,0.062);
    \node[font=\small,text=BookAmber,anchor=south]
      at (axis cs:60,0.062) {均值};
    \node[font=\small,text=BookInk,anchor=south]
      at (axis cs:40,0.067) {工况一};
    \node[font=\small,text=BookInk,anchor=south]
      at (axis cs:80,0.067) {工况二};
  \end{axis}
\end{tikzpicture}

  \caption{同一个双峰预测分布可以产生均值为$60\,^\circ\mathrm{C}$的点预测。实线为两个等权高斯密度的混合，每个分量的标准差为$3\,^\circ\mathrm{C}$；虚线标出均值。点预测没有保留两个高概率区域的位置。}
  \label{fig:pgm-overview-predictive-density}
\end{figure}

点预测和分布预测并不冲突。得到预测分布以后，仍可根据任务选择均值、概率最大的结果或其他行动。区别在于，分布先保留不确定性，点预测再按照使用目的把它压缩成一个结果。

困难在于，多个变量的联合分布会迅速变大。若有$d$个二值变量，完整联合表包含$2^d$个状态，需要$2^d-1$个自由参数。变量增多以后，有限数据很难覆盖所有组合，计算一个变量的边缘概率也可能需要汇总大量状态。概率图模型并没有取消这些困难，而是试图利用可信的局部依赖减少需要表示和计算的内容。

最强的简化是假设所有变量相互独立，但这会删掉故障与报警之间真正有用的联系。更合适的做法是只删除在特定条件下已经多余的依赖。例如，设备是否报警可以依赖是否过热和传感器是否失灵；一旦这两个状态已经给定，故障状态不必再直接进入报警的局部概率。这样的局部假设正是图结构要记录的信息。

\section{图模型上的表示}
\label{sec:pgm-overview-representation}
\label{sec:pgm-overview-basic-concepts}

本章使用四个二值随机变量：$F$表示设备故障，$H$表示过热，$A$表示报警，$S$表示传感器失灵，值$1$表示相应事件发生，值$0$表示未发生。联合概率$p(f,h,a,s)$描述四个变量同时取给定值的可能性，其中小写字母表示对应随机变量的取值。

从联合分布中只保留一部分变量，需要把其余变量的可能取值汇总起来。对离散变量，这种\term{边缘化}（marginalization）通过求和完成；获得证据后重新比较各种可能性，则使用条件概率。例如
\begin{equation}
  p(f,a)=\sum_{h,s}p(f,h,a,s),
  \qquad
  p(f\mid a)=\frac{p(f,a)}{p(a)},\quad p(a)>0.
  \label{eq:pgm-overview-marginal-conditional}
\end{equation}
这里对$h,s$求和，表示不把未知状态硬填成某个值，而是保留它们的所有可能性。连续变量采用相应的积分。

如果无论$X$取什么值，$Y$的分布都不改变，就称$X$与$Y$独立。更常用的是\term{条件独立}（conditional independence）：给定$Z$以后，继续观察$Y$不再改变对$X$的判断，记为$X\perp Y\mid Z$。在离散情形，它可以写成
\begin{equation}
  p(x,y\mid z)=p(x\mid z)p(y\mid z),
  \qquad p(z)>0.
  \label{eq:pgm-overview-conditional-independence}
\end{equation}
条件独立不是说$X$与$Y$永远没有关系，而是说$Z$已经包含了当前问题中连接它们所需的信息。概率图模型正是利用这类局部假设压缩联合分布。

在设备例子中，假设故障$F$与传感器失灵$S$独立，过热$H$只直接依赖故障，报警$A$只直接依赖过热和传感器状态。联合分布便可以写成
\begin{equation}
  p(f,h,a,s)=p(f)p(s)p(h\mid f)p(a\mid h,s).
  \label{eq:pgm-overview-device-joint}
\end{equation}
只要每个局部条件分布本身合法，这四个函数的乘积就定义了完整联合分布。对二值变量，两个根变量各需一个参数，过热条件分布需要两个，报警条件分布需要四个，总共$8$个自由参数，少于完整联合表的$15$个。减少的参数来自局部依赖假设，而不是记号上的省略。

把每个随机变量画成节点，并用箭头表示“哪个变量进入哪个局部条件分布”，就得到图~\ref{fig:pgm-overview-device-graph}。例如$H\to A\leftarrow S$对应局部项$p(a\mid h,s)$。箭头在这里记录概率分解，不凭图形本身断言因果关系。

\begin{figure}[htbp]
  \centering
  \begin{tikzpicture}[
  x=1mm,y=1mm,
  every node/.style={font=\small,text=BookInk},
  pgm variable/.style={
    circle,draw=BookTeal,line width=0.8pt,
    fill=BookPaper,minimum size=10mm,inner sep=1pt
  }
]
  \node[pgm variable] (fault) at (0,0) {$F$};
  \node[pgm variable] (heat) at (27,0) {$H$};
  \node[pgm variable] (alarm) at (54,0) {$A$};
  \node[pgm variable] (sensor) at (81,0) {$S$};
  \draw[book arrow] (fault) -- (heat);
  \draw[book arrow] (heat) -- (alarm);
  \draw[book arrow] (sensor) -- (alarm);
  \node[anchor=north] at (0,-7) {故障};
  \node[anchor=north] at (27,-7) {过热};
  \node[anchor=north] at (54,-7) {报警};
  \node[anchor=north] at (81,-7) {传感器失灵};
  \node[anchor=south,text=BookMuted] at (40.5,10)
    {$F,S$为独立根变量};
\end{tikzpicture}

  \caption{设备诊断模型。节点表示随机变量，箭头指向局部条件分布所描述的变量，例如$H\to A\leftarrow S$对应$p(a\mid h,s)$。箭头指定概率分解，不自动具有因果含义。}
  \label{fig:pgm-overview-device-graph}
\end{figure}

图只说明局部依赖结构，还需要局部概率函数给出具体数值。反过来，一组数值表若没有说明各表对应哪些变量，也不足以解释整个模型。理解概率图模型时，应同时分清图结构和组成联合分布的局部函数：前者记录依赖怎样组织，后者决定模型实际给出什么概率。

图模型最基本的概率表示有两类。\term{有向图模型}（directed graphical model）像设备模型一样，用箭头组织局部条件分布；后续第~\ref{chap:pgm-directed}章将说明怎样从有向图读取条件独立关系。\term{无向图模型}（undirected graphical model）用无向边表示对称的局部相互作用，再将一组非负局部函数的乘积归一化为概率分布；第~\ref{chap:pgm-undirected}章将展开这种表示。

\term{因子图}（factor graph）关注的是局部乘积本身。它把随机变量和局部函数画成两类节点，并在变量属于某个函数的输入时连边。因子图可以表示有向模型的条件概率乘积，也可以表示无向模型的局部函数乘积，因此它是组织计算的表示方式，不是与有向图、无向图互斥的第三类概率语义。第~\ref{chap:pgm-representations}章将比较三种表示保留的信息。

\section{图模型上的推断与学习}
\label{sec:pgm-overview-inference-learning}

图结构和局部概率确定以后，模型仍需经过计算才能回答具体问题。\term{概率推断}（probabilistic inference）是在模型已知时，根据已有证据计算所需概率。要回答“报警后设备故障的概率是多少”，需要计算$p(F=1\mid A=1)$；由于过热$H$和传感器状态$S$没有被观测，计算时必须汇总它们的所有可能取值。

把设备模型扩展到更多部件。扩展后共有$30$个未观测的二值状态，其中包括原来的过热$H$、传感器状态$S$和新增的部件状态或运行状态；以下将这$30$个未观测状态统一记为$\bm{Z}=(Z_1,\ldots,Z_{30})$。报警$A=1$仍是证据，故障$F$仍是查询目标。后验需要把这些变量的每一种配置都纳入比较：
\begin{equation}
  p(F=1\mid A=1)
  =
  \frac{
    \displaystyle\sum_{\bm{z}\in\{0,1\}^{30}}
    p(F=1,A=1,\bm{z})
  }{
    \displaystyle\sum_{f\in\{0,1\}}
    \sum_{\bm{z}\in\{0,1\}^{30}}
    p(f,A=1,\bm{z})
  }.
  \label{eq:pgm-overview-expanded-diagnosis}
\end{equation}
若直接展开，分子要汇总$2^{30}$个未观测状态配置，约为$10.7$亿项；分母还要同时遍历故障的两个取值，共有$2^{31}$个完整配置项，约为$21.5$亿项。这里的项数表示求和中需要纳入多少种状态组合，不是模型必须具有同样多的参数。

局部因子使单个完整配置$p(f,A=1,\bm{z})$通常容易评价：只需查出或计算该配置涉及的局部概率，再将它们相乘。但后验不是某个完整配置的概率，而是对大量未观测状态配置的汇总。因而，“单个完整配置易评价”和“汇总全部未观测状态仍昂贵”可以同时成立。图中的局部结构有时能够复用计算，但能节省多少取决于依赖怎样连接；具体算法与结构条件留到后续章节。

\term{精确推断}（exact inference）要求在给定模型下得到原查询的准确答案。它适合结构允许有效计算、或问题规模较小的情形；结构复杂时，中间计算仍可能非常大。\term{近似推断}（approximate inference）则在无法承担精确计算时，接受一定误差，以换取时间或空间上的可行性。第~\ref{chap:pgm-elimination}章至第~\ref{chap:pgm-sampling}章将分别展开精确与近似推断，本章不提前介绍具体实现路线。

模型中的局部概率通常不是预先知道的，还需要从数据获得。\term{参数学习}（parameter learning）在图结构和局部函数形式已经确定时，估计这些函数中的数值。例如，若记录了许多设备的故障和过热状态，就可以估计$p(H=1\mid F=1)$。数据越少、状态越稀有，估计通常越不稳定；数据中存在缺失或未观测变量时，所需计数也无法直接得到。

如果连哪些变量应直接相连也不确定，就需要\term{结构学习}（structure learning）根据数据和领域约束选择图。结构过于简单可能遗漏重要依赖，结构过于复杂则通常需要更多数据并增加计算负担。学习到有向边也不自动等于发现因果方向；因果解释需要额外假设和证据。

推断与学习解决的是不同问题：推断固定模型、计算未知状态；学习固定数据、确定模型结构或参数。当数据含有未观测变量时，两者会互相依赖，因为学习参数之前，往往要先用当前模型推断这些变量可能取什么值。第~\ref{chap:pgm-parameter-learning}和第~\ref{chap:pgm-structure-learning}章将分别展开参数与结构学习。

\section{概率图模型的演进}
\label{sec:pgm-overview-history}

概率图模型并不是一套方法按年份依次替代另一套方法的结果。图~\ref{fig:pgm-overview-history}从现代视角出发，按照本书关心的问题将相关工作归纳为五条相互衔接、但至今仍并存的线索。这种组织方式用于说明不同模型和技术回应了哪些问题、留下了什么边界，并非逐项复原原作者的提出动机。

\begin{figure*}[htbp]
  \centering
  \begin{tikzpicture}[
  x=1mm,y=1mm,
  every node/.style={font=\footnotesize,text=BookInk},
  stage/.style={
    rectangle,
    draw=BookRule,
    fill=BookPaper,
    minimum width=29mm,
    minimum height=38mm,
    text width=25mm,
    align=center,
    inner sep=2mm
  },
  stage arrow/.style={
    -{Stealth[length=2mm]},
    draw=BookTeal,
    line width=0.9pt
  }
]
  \node[stage] (local) at (15,0) {
    \textcolor{BookTeal}{\heiti 局部随机过程}\\[1mm]
    {\scriptsize 回顾线索：整体状态过大}\\
    {\scriptsize 技术：时间与邻域局部性}\\
    {\scriptsize 边界：缺少统一图语义}
  };
  \node[stage] (semantics) at (48,0) {
    \textcolor{BookTeal}{\heiti 图结构语义}\\[1mm]
    {\scriptsize 回顾线索：依赖难以检查}\\
    {\scriptsize 技术：分离与局部分解}\\
    {\scriptsize 边界：图不自动等于因果}
  };
  \node[stage] (exact) at (81,0) {
    \textcolor{BookTeal}{\heiti 精确推断}\\[1mm]
    {\scriptsize 回顾线索：全局查询昂贵}\\
    {\scriptsize 技术：变量消除与团树}\\
    {\scriptsize 边界：树宽控制指数代价}
  };
  \node[stage] (learning) at (114,0) {
    \textcolor{BookTeal}{\heiti 学习与近似推断}\\[1mm]
    {\scriptsize 回顾线索：模型未知或难算}\\
    {\scriptsize 技术：参数、结构与近似}\\
    {\scriptsize 边界：近似和估计有误差}
  };
  \node[stage] (neural) at (147,0) {
    \textcolor{BookTeal}{\heiti 神经参数化}\\[1mm]
    {\scriptsize 回顾线索：局部函数受限}\\
    {\scriptsize 技术：VAE、规范化流与扩散}\\
    {\scriptsize 边界：归一化、似然与推断代价}
  };

  \draw[stage arrow] (local.east) -- (semantics.west);
  \draw[stage arrow] (semantics.east) -- (exact.west);
  \draw[stage arrow] (exact.east) -- (learning.west);
  \draw[stage arrow] (learning.east) -- (neural.west);
\end{tikzpicture}

  \caption{按本书问题线索组织的概率图模型现代回顾。图中的“回顾线索”联系问题、技术与边界；箭头表示讨论的问题逐步扩展，不表示后出现的技术取代了此前方法。}
  \label{fig:pgm-overview-history}
\end{figure*}

当完整轨迹或整体状态空间很大时，可以用局部规律描述相邻时刻或相邻位置的依赖。沿着\term{局部随机过程}（local stochastic process）这条线索，Markov性质用当前状态概括与未来有关的历史，格点模型则用邻域相互作用描述整体配置；Ising模型是后一类模型的早期代表。局部性减少了必须直接指定的关系，却还没有给出后来概率图模型中统一的图分离语义，也不保证所有全局概率都容易计算\scite{pgm-overview-ising1925}

第二条线索是怎样明确表示“哪些变量直接相关、给定什么以后哪些信息变得多余”。\term{图结构语义}（graph-structured semantics）把条件独立、局部概率分解和图上的分离规则联系起来；有向网络与无向网络由此成为组织复杂联合分布的两种基本方式。Pearl对有向概率网络的系统化处理以及Lauritzen对图模型理论的整理，确立了这一表示框架。图能编码概率约束，但有向边本身不自动获得因果含义，特殊参数也可能产生图中没有表达的额外独立性\scite{pgm-pearl1988,pgm-lauritzen1996}

第三条线索关注怎样由局部函数得到边缘概率和后验。\term{精确推断}（exact inference）中，变量消除通过重排求和减少重复计算，树上的置信传播和一般图上的团树则把中间结果组织成可复用的局部消息；Lauritzen与Spiegelhalter将图变换和局部概率计算结合为适用于一般网络的系统方法。精确推断没有消除组合爆炸：对离散表格而言，中间团的状态空间随树宽指数增长，高树宽模型仍可能无法承受\scite{pgm-overview-lauritzen1988,pgm-koller2009}

第四条线索关注模型参数或图结构未知、记录不完整、精确后验不可承受时，怎样从数据估计模型并在预算内近似查询。\term{学习与近似推断}（learning and approximate inference）中，参数学习估计局部函数，结构学习选择依赖连接；存在缺失或潜变量时，期望最大化（expectation-maximization, EM）等方法交替计算后验统计量和更新参数，变分方法和采样方法则以优化误差或随机误差换取可计算性。这些路线分别受到组合搜索与可识别性、局部最优、近似族或有限样本波动的限制，不是对精确推断的无条件替代\scite{pgm-overview-dempster1977,pgm-wainwright2008,pgm-koller2009}

第五条线索关注怎样表达图像、文本等高维数据中的复杂局部关系。\term{神经参数化}（neural parameterization）用神经网络表示局部条件分布、势函数或近似后验；VAE将神经生成模型与可学习的近似后验结合，规范化流用可逆神经变换构造可评价密度或增强近似后验，扩散模型则学习逐步逆转人为噪声过程。表达能力提高以后，归一化、似然评价以及推断或采样代价仍需分别处理\scite{kingma2013,pgm-overview-rezende2015,pgm-overview-sohl2015}

\section{本章小结}
\label{sec:pgm-overview-summary}

点预测给出一个结果，概率分布则保留其他结果的可能性。多个变量的联合分布会随变量数量迅速增大，概率图模型因此用局部条件独立组织依赖，再由局部概率函数组成联合分布。图结构减少了需要明确指定的关系，但这种简化是否成立，取决于局部假设是否符合问题。

图模型上的表示以有向图和无向图为两种基本概率语义，因子图则把两类模型中的局部乘积显式画出来。局部因子可以让单个完整配置易于评价，却不保证汇总大量未观测状态配置同样容易。模型给定后，推断根据证据回答概率查询；数据给定后，学习估计参数或选择结构。

表示、计算和学习上的困难共同推动了概率图模型的演进：局部随机规律压缩依赖，图结构赋予这些依赖可检查的语义，精确推断重组全局计算，学习与近似推断处理未知模型和不可承受的查询，神经参数化再扩展局部函数的表达能力。这些路线各有边界，没有按时代相互取代。下一章将从有向图出发，严格说明图结构、条件独立和联合分解怎样联系起来。
